% R008 — Guide to Cultivating Complex Analysis (Bahasa Indonesia)
% Reader-facing translation unit: source/ca-v1.9/ca.tex lines 1162–1175.
% Source release: v1.9 (commit a4d25d9ba37a486a5a020219eb8bd4e6602fd14c).
% Translation provenance: OpenAI Codex gpt-5.6-sol, Ultra, acting on the user's request.
% Preserve source notation and paragraph structure.

Representasi bilangan kompleks ini cukup sering muncul dalam penerapan.
Sebagai contoh,
matriks kompleks berukuran $m \times n$ dapat direpresentasikan oleh
matriks real berukuran $2m \times 2n$ dengan mengganti setiap entri dengan
matriks berukuran $2 \times 2$.
Dengan demikian, perangkat lunak yang disiapkan untuk bekerja dengan matriks
real dapat dengan mudah dikelabui agar bekerja dengan matriks kompleks.

Bagi kita, penerapan utamanya adalah memahami turunan suatu
fungsi bernilai kompleks dari variabel kompleks $f \colon \C \to \C$.
Dengan memandang fungsi tersebut sebagai pemetaan $f \colon \R^2 \to \R^2$,
turunan real $f$ merupakan matriks $2 \times 2$.
Objek kajian analisis kompleks, yaitu fungsi holomorf (atau analitik),
adalah fungsi-fungsi yang matriks turunannya yang real bersesuaian dengan
perkalian oleh suatu bilangan kompleks.
